\documentclass[10pt,a4paper]{amsart}
\usepackage[margin=19mm]{geometry}
\usepackage{amsmath,amssymb,mathtools}
\usepackage[scr=rsfs,cal=euler]{mathalfa}
\usepackage{xfrac}
\usepackage[unicode,hidelinks,bookmarks=false]{hyperref}
\newcommand{\ie}{i.e.\ }
\newcommand{\cf}{cf.\ }
\newcommand{\resp}{respectively\ }
\newcommand{\Subsection}{\subsection}
\providecommand{\nobreakdash}{\nobreak\hbox{-}}
\setlength{\emergencystretch}{2em}
\multlinegap0pt
\usepackage{bm}
\usepackage{bbm}
\usepackage{dsfont}
\usepackage{xspace}
\usepackage{cancel}
\usepackage{mathtools}
\usepackage{amsthm}
\makeatletter
\@ifundefined{lem}{\newtheorem{lem}{Lemma}}{}
\@ifundefined{prop}{\newtheorem{prop}{Proposition}}{}
\makeatother
\newcommand{\Ds}{(-\Delta)^{s}}
\newcommand{\R}{\mathbb{R}}
\newcommand{\Tss}{\mathbf{T}^{(\sigma)}_\lambda}
\newcommand{\Ts}{\mathbf{T}}
\newcommand{\be}{\begin{equation}}
\newcommand{\ee}{\end{equation}}
\newcommand{\eps}{\varepsilon}
\newcommand{\eu}{e}
\newcommand{\pt}{\partial}
\renewcommand{\geq}{\geqslant}
\renewcommand{\leq}{\leqslant}
\title[Existence and Uniqueness for the Fractional Gelfand Equation]{Existence and Uniqueness for the Fractional Gelfand Equation in $\mathbb{R}$}
\author{Florian P. Lanz, Enno Lenzmann}
\date{Source: arXiv:2506.07577v1 (CC BY 4.0)}
\begin{document}
\maketitle

\noindent\emph{Source and licence.} Existence and Uniqueness for the Fractional Gelfand Equation in $\mathbb{R}$. Version arXiv:2506.07577v1 (CC BY 4.0), distributed under the Creative Commons Attribution 4.0 International licence, as recorded in the arXiv licence metadata for this exact version and shown on the abstract page https://arxiv.org/abs/2506.07577v1. Full rights notice: licenses/arxiv-2506.07577-CC-BY-4.0.txt.\par\smallskip

\noindent\textit{Editorial scope.} Counted unit: the Lem whose source label is \texttt{lem:nice} in the article (source0.tex lines 992--994, source label \texttt{lem:nice}), delivered together with the article's own complete original proof of it (source0.tex lines 996--1070, one \texttt{proof} environment of 75 lines). No mathematics is written, repaired or re-derived: statement and proof are the article's own text line for line. Required context retained uncounted: lem:decay (lines 645--651); prop:T\_bounds (lines 795--804); lem:HLS (lines 924--930); lem:compact (lines 1486--1492); lem:helly (lines 1498--1507). Every definition, display and label of those lines is retained unchanged. Omitted from this package and not redistributed: 1--644, 652--794, 805--923, 931--991, 995--995, 1071--1485, 1493--1497, 1508--1795. Nothing inside a retained excerpt is omitted or altered: every retained body is delivered exactly as the article's lines are written, so no omission receipt of any kind accompanies this package. Citations: the retained excerpts carry no citation command at all, so no bibliography entry of the article is transcribed and no reference is redistributed. Reference adaptation: none was needed and none was made; every reference inside the delivered unit resolves inside it, so no unresolved reference or citation is produced. Printed slips kept literal: no printed word, symbol, hypothesis or number of the counted statement or of its proof was altered. Layout and class: the article's own class is \texttt{amsart} with packages including amsfonts, amsmath, amssymb, amsthm, color, bbm, nccmath, mathrsfs, dsfont, babel, esint, graphicx, hyperref, geometry; the wrapper is the lane's pinned \texttt{amsart} editorial wrapper at 10pt on A4, which restates the theorem-like environments the retained text needs and adds the article's own macro definitions that the retained text needs (\texttt{\textbackslash Ds}, \texttt{\textbackslash R}, \texttt{\textbackslash Ts}, \texttt{\textbackslash Tss}, \texttt{\textbackslash be}, \texttt{\textbackslash ee}, \texttt{\textbackslash eps}, \texttt{\textbackslash eu}, \texttt{\textbackslash geq}, \texttt{\textbackslash leq}, \texttt{\textbackslash pt}), with extra wrapper packages (bm, bbm, dsfont, xspace, cancel, mathtools). The delivered numbering is the unit's own. Assets: no retained span contains a figure environment, a graphics-inclusion command, a PDF-page inclusion, a table or a TikZ picture; no image file is copied into this package, the package declares no asset dependency and no image file is redistributed with it. Licence: version arXiv:2506.07577v1 of this article is distributed under the Creative Commons Attribution 4.0 International licence, as recorded in the arXiv licence metadata for this exact version and shown on the abstract page https://arxiv.org/abs/2506.07577v1. A full rights notice is delivered at licenses/arxiv-2506.07577-CC-BY-4.0.txt. Characters: every retained excerpt is pure ASCII, so the delivered engine, XeLaTeX with its default Latin Modern text font, reports no missing character.
\begin{lem} \label{lem:decay}
Let $s \in (\frac 1 2, 1)$ and $\rho \in L^1(\R)$ with $\rho \geq 0$. Suppose $u \in L_s(\R)$ is an integral solution of $\Ds u = \rho$. Then, for any $\eps \in (0,1)$, there exists $R > 0$ such that
$$
u(x) \leq -c_\alpha (1-\eps) |x|^{2 \alpha} \left ( \int_\R \rho(y) \, dy \right ) + A \quad \mbox{for} \quad  |x| \geq R
$$
with some constant $A \in \R$.
\end{lem}

\begin{prop} \label{prop:T_bounds}
For all $v \in X_\alpha$, it holds
$$
0 < \Tss[v](x) \lesssim \lambda e^{-\frac{1}{2} \sigma^2 x^2} e^{-d_\alpha \| v \|_{L^2}^2 |x|^{2 \alpha} + 2 d_\alpha \| |x|^{\alpha} v \|_{L^2}^2} \quad \mbox{for $x \in \R$}
$$
with some constant $d_\alpha > 0$ depending only on $\alpha$.  Moreover, the function $x \mapsto \Tss[v](x)$ is even and $C^1$ with the pointwise bound
$$
|\pt_x \Tss[v](x)| \lesssim \lambda e^{-\frac{1}{2} \sigma^2 x^2} e^{-d_\alpha \| v \|_{L^2}^2 |x|^{2 \alpha} + 2 d_\alpha \| |x|^{\alpha} v \|_{L^2}^2} \left ( \sigma^2 |x| + \| v \|_{X^\alpha}^2 + 1 \right ) \quad \mbox{for $x \in \R$}.
$$
\end{prop}

\begin{lem} \label{lem:HLS}
Every fixed point $v \in X_\alpha$ of $\Tss$ satisfies the bound
$$
\| v \|_{L^p} \leq C(\lambda, p, \| K \|_{L^\infty}) \quad \mbox{for} \quad \mbox{$\frac{2}{1+ \alpha} \leq p \leq \infty$},
$$
with some constant $C=C(\lambda,p,\| K \|_{L^\infty}) > 0$ independent of $\sigma \in \R$.
\end{lem}

\begin{lem} \label{lem:compact}
Let $C > 0$ and $d >0$ be constants. Suppose that $\mathcal{F} \subset X_\alpha$ satisfies
$$
|v(x)| \leq C \eu^{-d |x|^{2 \alpha}} \quad |\pt_x v(x)| \leq C \quad \mbox{for all $v \in \mathcal{F}$ and a.\,e.~$x \in \R$} \, .
$$
Then $\mathcal{F}$ is relatively compact in $X_\alpha$. 
\end{lem}

\begin{lem} \label{lem:helly}
Let $1 \leq p_1 < p_2 < \infty$ be given. Suppose that $f_n=f_n^*$ is a sequence of symmetric-decreasing functions satisying
$$
\| f_n \|_{L^p} \leq C \quad \mbox{for all $p \in [p_1, p_2]$}
$$ 
with some constant $C>0$. Then there exists a subsequence $(f_{n_k})_{k=1}^\infty$ and a symmetric-decreasing function $f=f^*$ such that $f_{n_k} \to f$ almost everywhere in $\R$ and
$$
\mbox{$f_{n_k} \to f$ in $L^p(\R)$ for all $p \in (p_1, p_2)$}.
$$
\end{lem}

\begin{lem} \label{lem:nice}
Let $\lambda > 0$ be given. Suppose that $\sigma_n \neq 0$ is a sequence such that $\sigma_n \to 0$ and let $(v_n)_{n \geq 1} \subset X_\alpha$ be a sequence of fixed points of $\Ts_{\lambda}^{(\sigma_n)}$. Then, after passing to a subsequence, it holds that $v_{n} \to v$ in $X_\alpha$ and $v \not \equiv 0$ is a fixed point of $\Ts_{\lambda}^{(0)}$.
\end{lem}

\begin{proof}
We divide the proof into the following steps.

\medskip
\textbf{Step 1.} From the uniform bounds $\| v_n \|_{L^p}$ for $p \in [\frac{2}{1+\alpha}, \infty]$ derived in Lemma \ref{lem:HLS} and the fact that the $v_n = v_n^*$ are symmetric-decreasing functions, we infer from Lemma \ref{lem:helly} (after passing to a subsequence) that
\be
\mbox{$v_n \to v$ in $L^p(\R)$ for any $\frac{2}{1+\alpha} < p < \infty$ and $v_n \to v$ almost everywhere in $\R$}
\ee
with some symmetric-decreasing function $v=v^*$ with $0 \leq v(x) \leq \lambda \sqrt{K(0)}$ for a.\,e.~$x \in \R$. We claim that
$$
e^{-\frac{1}{2} \int_0^x H_\alpha(v_n^2)(y) dy }\to e^{-\frac{1}{2} \int_0^x H_\alpha(v^2)(y) \,dy } \quad \mbox{for almost every $x \in \R$}.
$$ 
To see this, we note that $\| H_\alpha(v_n^2) \|_{L^\infty} \lesssim \| v_n^2 \|_{L^1 \cap L^\infty} = \| v_n \|_{L^2 \cap L^\infty} \lesssim 1$. Thus, by dominated convergence,
$$
\lim_{n \to \infty} \int_0^x H_\alpha(v_n^2)(y) dy = \int_0^x \lim_{n \to \infty} H_\alpha(v_n^2)(y) dy = \int_0^x H_\alpha(v^2)(y) dy
$$
for every $x \in \R$, where in the second step we use that $H_\alpha(v_n^2)(y)\to H_\alpha(v^2)(y)$ for a.\,e.~$y \in \R$, which follows from $v_n \to v$ in $L^p$ for any $p \in (2, \infty)$ and the properties of $H_\alpha$.

In the equation
\be \label{eq:vn}
v_n(x) = \lambda \sqrt{K(x)} e^{-\frac{1}{2} \sigma_n^2 x^2} e^{-\frac{1}{2} \int_0^x H_\alpha(v_n^2)(y) dy}
\ee
we can pass to the limit $n \to \infty$ on both sides of the equation to conclude that the limit $v=v^* \in L^2(\R) \cap L^\infty(\R)$ satisfies
\be \label{eq:limit}
v(x) = \lambda \sqrt{K(x)} e^{-\frac{1}{2} \int_0^x H_\alpha(v^2)(y) dy} \quad \mbox{for almost every $x \in \R$}.
\ee
In particular, this equation shows that $v \not \equiv 0$.
 
 \medskip
 \textbf{Step 2.} We demonstrate that there exist constants $C, d > 0$ such that
 \be \label{eq:v_n_bounds}
 |v_n(x)| \leq  C e^{-d |x|^{2 \alpha}}, \quad |\pt_x v_n(x)| \leq C \quad \mbox{for all $x \in \R$ and $n \geq 1$}.
 \ee
To this end, we consider the sequence of functions
$$
w_n(x) = -\frac{1}{2} \int_0^x H_\alpha(v_n^2)(y) \, dy = -c_\alpha \int_\R (|x-y|^{2 \alpha} - |y|^{2 \alpha}) v_n(y)^2 \, dy .
$$
Since $v_n \to v$ in $L^2(\R)$ with $v \not \equiv 0$, there exists a constant $c>0$ and $n_0 \geq 1$ so that
$$
\int_{\R} v_n(y)^2 \, dy \geq c \quad \mbox{for all $n \geq n_0$}.
$$
Moreover, by the uniform decay estimate $v_n(x)^2 \lesssim |x|^{-2/p}$ with $p=\frac{2}{1+\alpha} < 2$, we find that the sequence $(v_n)$ is {\em tight} in $L^2(\R)$, i.\,e., for every $\eps > 0$ there exists $M > 0$ such that
$$
\int_{|y| > M} v_n(y)^2 \, dy \leq \eps \quad \mbox{for all $n \geq 1$}.
$$ 
Next, we argue as in the proof of Lemma \ref{lem:decay}. We can take some $M>0$ such that
$$
\int_{|y| > M} v_n(y)^2 \, dy \leq \frac{c}{4} \quad \mbox{and} \quad \int_{|y| \leq M} v_n(y)^2 \, dy \geq \frac{3 c}{4} \quad \mbox{for all $n \geq n_0$}.
$$
Let $\delta > 0$ be sufficiently small such that $|1-t|^{2 \alpha} - |t|^{2 \alpha} \geq \frac{1}{2}$ for $|t| \leq \delta$. Thus if we define $R:=\delta^{-1} M > 0$, we obtain
$$
|x-y|^{2 \alpha} - |y|^{2 \alpha} \geq \frac{1}{2} |x|^{2 \alpha} \quad \mbox{for $|y| \leq M$ and $|x| \geq R$}.
$$
Using that $v_n^2 \geq 0$ and the general bound $|x-y|^{2\alpha} -|y|^{2 \alpha} \geq -|x|^{2 \alpha}$, we deduce
\begin{align*}
& \int_\R (|x-y|^{2 \alpha}- |y|^{2 \alpha}) v_n(y)^2 \, dy \geq \frac{1}{2} |x|^{2 \alpha} \int_{|y| \leq M} v_n(y)^2 \, dy - |x|^{2 \alpha} \int_{|y| > M} v_n(y)^2 \, dy \\
& \geq \frac{3 c}{8} |x|^{2 \alpha} - \frac{c}{4} |x|^{2 \alpha} = \frac{c}{8} |x|^{2 \alpha} \quad \mbox{for $|x| \geq R$}.
\end{align*}
Choosing now $d = c_\alpha \cdot \frac{c}{8} >0$, we obtain the bound
$$
w_n(x) \leq -d |x|^{2 \alpha} \quad \mbox{for $|x| \geq R$ and $n \geq n_0$}.
$$
Therefore,
$$
0 < v_n(x) = \lambda \sqrt{K(x)} e^{-\frac{1}{2} \sigma_n^2 x^2} e^{w_n(x)}  \leq \lambda \| K \|_{L^\infty}^{1/2}  e^{-d |x|^{2 \alpha}} \quad \mbox{for $|x| \geq R$ and $n \geq n_0$}.
$$

Since also $0 < v_n(x) \leq v_n(0) = \lambda \sqrt{K(0)} = \lambda \| K \|_{L^\infty}^{1/2}$ for all $x \in \R$, we readily deduce that the first estimate in \eqref{eq:v_n_bounds} holds for some sufficiently large constant $C>0$ independent of $n$. Also, from the general estimate in Proposition \ref{prop:T_bounds}, we deduce
$$
|\pt_x v_n(x)| \leq C e^{-d|x|^{2 \alpha}} (\sigma_n^2 |x| + 1) \leq C \quad \mbox{for all $x \in \R$ and $n \geq 1$}
$$
with some constant $C >0$ independent of $n$. This proves \eqref{eq:v_n_bounds}.

By Lemma \ref{lem:compact}, we infer (after passing to a subsequence) that $v_n \to v$ in $X_\alpha$ and $v \not \equiv 0$ is a fixed point of $\Ts_\lambda=\Ts_{\lambda}^{(0)}$ thanks to \eqref{eq:limit}. The proof of Lemma \ref{lem:nice} is now complete.
\end{proof}

\end{document}
